% Part: first-order-logic
% Chapter: syntax-and-semantics
% Section: main-operator

\documentclass[../../../include/open-logic-section]{subfiles}

\begin{document}

\olfileid{fol}{syn}{mai}

\olsection{সূত্রের \printtoken{S}{main operator}}

\begin{explain}
!!a{formula}~$!A$ নির্মাণে সব শেষে ব্যবহৃত অপারেটরটি নিয়ে কথা বলা
প্রায়ই উপযোগী। এই অপারেটরকে $!A$-এর \emph{প্রধান অপারেটর} বলে।
স্বজ্ঞাতভাবে, এটি $!A$-এর “সবচেয়ে বাইরের” অপারেটর। যেমন,
$\lnot !A$-এর প্রধান অপারেটর $\lnot$, $(!A \lor !B)$-এর প্রধান
অপারেটর $\lor$, ইত্যাদি।
\end{explain}


\begin{defn}[!!^{main operator}]
\ollabel{def:main-op}
!!a{formula}~$!A$-এর \emph{!!{main operator}} নিচের নিয়মে সংজ্ঞায়িত:
\begin{enumerate}
\item \indcase*{!A}{!A}{$\indfrm$-এর কোনো !!{main operator} নেই।}

\tagitem{prvNot}{\indcase{!A}{\lnot !B}{$\indfrm$-এর !!{main operator}
  হলো~$\lnot$।}}{}

\tagitem{prvAnd}{\indcase{!A}{(!B \land !C)}{$\indfrm$-এর
  !!{main operator} হলো~$\land$।}}{}

\tagitem{prvOr}{\indcase{!A}{(!B \lor !C)}{$\indfrm$-এর
  !!{main operator} হলো~$\lor$।}}{}

\tagitem{prvIf}{\indcase{!A}{(!B \lif !C)}{$\indfrm$-এর
  !!{main operator} হলো~$\lif$।}}{}

\tagitem{prvIff}{\indcase{!A}{(!B \liff !C)}{$\indfrm$-এর
  !!{main operator} হলো~$\liff$।}}{}

\tagitem{prvAll}{\indcase{!A}{\lforall[x][!B]}{$\indfrm$-এর
  !!{main operator} হলো~$\lforall$।}}{}

\tagitem{prvEx}{\indcase{!A}{\lexists[x][!B]}{$\indfrm$-এর
  !!{main operator} হলো~$\lexists$।}}{}
\end{enumerate}
\end{defn}

প্রত্যেক ক্ষেত্রে আমরা সূত্রে !!{main operator}-টির নির্দিষ্ট নির্দেশিত
\emph{সংঘটন} বোঝাই। যেমন, $((!D \lif !E) \lif (!E \lif !D))$ সূত্রটি
$(!B \lif !C)$ রূপের, যেখানে $!B$ হলো $(!D \lif !E)$ এবং $!C$ হলো
$(!E \lif !D)$; তাই $\lif$-এর দ্বিতীয় সংঘটনটি !!{main operator}।

\begin{explain}
এটি এমন একটি ফলনের \emph{পুনরাবৃত্তিমূলক} সংজ্ঞা, যা প্রত্যেক
অপরমাণু !!{formula}s-কে তার !!{main operator}-এর সংঘটনে প্রতিচিত্রিত
করে। !!{formula}s-কে যেভাবে আরোহীভাবে সংজ্ঞায়িত করা হয়েছে, তার ফলে
প্রত্যেকটি !!{formula}~$!A$ \olref{def:main-op}-এর কোনো একটি ক্ষেত্র
পূরণ করে। এতে প্রত্যেক অপরমাণু !!{formula}~$!A$-এর জন্য !!a{main
  operator}-এর অস্তিত্ব নিশ্চিত হয়। প্রত্যেকটি !!{formula} এই
শর্তগুলির কেবল একটি পূরণ করে এবং প্রতিটি ক্ষেত্রে যে ছোট
!!{formula}s থেকে $!A$ নির্মিত সেগুলি এককভাবে নির্ধারিত; তাই $!A$-এ
!!{main
  operator}-টির সংঘটনও একক এবং আমরা সত্যিই একটি ফলন সংজ্ঞায়িত
করেছি।
\end{explain}

কোন সংকেতটি !!{main operator}, তার ভিত্তিতে আমরা !!{formula}s-কে
\olref{tab:main-op}-এর নামগুলি দিই।
\iftag{defNot,defOr,defAnd,defIf,defIff,defTrue,defFalse,defEx,defAll}
{ তবে মনে রাখবেন, সংজ্ঞায়িত অপারেটর বিধিবদ্ধভাবে !!{formula}s-এ
দেখা যায় না। সেগুলি কেবল সংক্ষিপ্ত রূপ; তাই বিধিবদ্ধভাবে তারা কোনো
সূত্রের প্রধান অপারেটর হতে পারে না। ফলে সব !!{formula}s-সংক্রান্ত
প্রমাণে সেগুলিকে আলাদাভাবে বিবেচনা করার দরকারও হয় না।}

\begin{table}[!h]
\centering
\begin{tabular}{c | c | c}
!!^{main operator} & !!{formula}-এর প্রকার & উদাহরণ\\
\hline
কোনোটিই নয় & পরমাণু (!!{formula}) &
\iftag{prvFalse}{$\lfalse$,}{}
\iftag{prvTrue}{$\ltrue$,}{}
$\Atom{R}{t_1, \dots, t_n}$,
$\eq[t_1][t_2]$\\
$\lnot$ & নঞর্থকরণ & $\lnot !A$ \\
$\land$ & সংযোজন & $(!A \land !B)$ \\
$\lor$ & বিয়োজন & $(!A \lor !B$) \\
$\lif$ & !!{conditional} & $(!A \lif !B$) \\
$\liff$ & !!{biconditional} & $(!A \liff !B)$ \\
$\lforall[][]$ & সার্বিক (!!{formula})& $\lforall[x][!A]$ \\
$\lexists[][]$ & অস্তিত্বসূচক (!!{formula})& $\lexists[x][!A]$
\end{tabular}
\caption{!!{formula}s-এর প্রধান অপারেটর ও নাম}
\ollabel{tab:main-op}
\end{table}
% BN-SRC-089 TARGET-CORRECTION-BEGIN
\par\smallskip\noindent\textnormal{\small\textbf{উৎস-সংশোধন (BN-SRC-089):} হিমায়িত ইংরেজি উৎসের সংযোজন-সারির উদাহরণে গণিত-মোডটি ডান বন্ধনীর আগে শেষ হয়েছে: $(!A \land !B$)। বাংলা পাঠে ডান বন্ধনীটি একই গণিত-রাশির ভিতরে রেখে সুগঠিত সূত্র $(!A \land !B)$ করা হয়েছে।} % BN-SRC-089 TARGET-CORRECTION-END

\end{document}
